% Additive, independent reconstruction of incoming LRC6--10.
% Original WCF/NCT/TAC files are unchanged. Human bibliography keys retained.
\begingroup
\section{Uniform all-rank transport for the original conductor}
\label{sec:LRC-transport-audit}

Fix the original simple quartet, period, branch and conductor coefficients.
The programme origin retains its original attribution
\cite{human:globalization2026,human:splitzero2026}. Put
\[
\begin{gathered}
 a=4\ell+1\ge9,\qquad 0<\delta<\tfrac12,\qquad \gamma>2,
 \qquad q=(a+1)^2,\qquad Q=(a-7)^2,\\
 \beta_{n,u,w}=\tfrac n2+(2u-n)\delta+i(2w-n)\gamma,
 \qquad E_A(z)=\sum_{u,w=0}^{8}a_{uw}e^{\beta_{8,u,w}z},\\
 \mu_j=E_A^{(j)}(0),\qquad
 v=\operatorname{ord}_0E_A\le80,\quad \mu_v\ne0,
 \quad A_{\rm abs}=\sum_{u,w=0}^{8}|a_{uw}|,\\
 s\in\{1,a\},\quad b=s/2,\quad c_0=a/2,\quad c_1=a/2-4,
 \quad D\in\mathbb Z_{\ge0},\quad N=D+v,\quad d=D+1.
\end{gathered}                                                   \tag{LT1}
\]
The actual coefficients, $v,\delta,\gamma$ are fixed as $a$ varies.
The inner product is conjugate-linear in its first argument. Retain
the complete original measure and all its factors:
\[
\begin{gathered}
 dm_s(y)=\frac{(2\pi)^{s/2}2^{s/2}}{4\pi\Gamma(s/2)}
       |\Gamma(s/4+iy/2)|^2\,dy,\qquad M_s=(2\pi)^{s/2},\\
 \|P\|_{s,c}^2=\int_{\mathbb R}|P(c+iy)|^2\,dm_s(y),\qquad
 p_{0,s}=1,\quad p_{1,s}=y,\\
 p_{n+1,s}=yp_{n,s}-n(n-1+s/2)p_{n-1,s},\qquad
 h_{n,s}=M_sn!(s/2)_n,\\
 \phi_{n,s,c}(S)=p_{n,s}((S-c)/i)/\sqrt{h_{n,s}},\qquad
 [S^n]\phi_{n,s,c}=i^{-n}/\sqrt{h_{n,s}},\qquad
 \rho_n=\sqrt{n!/(s/2)_n}.
\end{gathered}                                                   \tag{LT2}
\]
These are the original orthonormal frames. To check the classical
dictionary with its mass, take
$p_{n,s}(y)=n!P_n^{(s/4)}(y/2;\pi/2)$ in the Meixner--Pollaczek
orthogonality formulas, DLMF (18.19.8)--(18.19.9)
\cite{human:dlmf18}. Their hypotheses hold because $s/4>0$.
With $y=2x$, the squared norm is exactly
\[
 \frac{M_s2^{s/2}}{4\pi\Gamma(s/2)}\,
 2(n!)^2\frac{2\pi\Gamma(n+s/2)}{2^{s/2}n!}
 =M_sn!(s/2)_n .
\]
The generating identity at these same parameters is DLMF (18.23.7)
\cite{human:dlmf18}:
\[
 \sum_{n\ge0}\frac{p_{n,s}(y)}{n!}t^n
 =(1+t^2)^{-s/4}e^{y\arctan t},\qquad |t|<1.          \tag{LT3}
\]
It also follows directly by multiplying the recurrence in LT2 by
$t^n/n!$: the series satisfies
$(1+t^2)\partial_t F=(y-(s/2)t)F$, $F(y,0)=1$.
The analytic factors and $\arctan$ take their branches fixed at zero.

For $\mathcal P_j=\mathbb C[S]_{\le j}$ and
$\mathcal P_{-1}=0$, the literal conductor is
\[
 (\mathcal T_AP)(S')=\sum_{u,w=0}^{8}a_{uw}P(S'+\beta_{8,u,w}),
 \qquad
 A_{D,s}:(\mathcal P_N/\mathcal P_{v-1},\|\cdot\|_{s,c_0,\rm quot})
       \longrightarrow(\mathbb C[S']_{\le D},\|\cdot\|_{s,c_1}).
                                                               \tag{LT4}
\]
Indeed $\mathcal T_AS^n=\sum_{j=v}^n\binom nj\mu_jS'^{n-j}$.
The leading coefficient for $n\ge v$ is $\binom nv\mu_v\ne0$,
and all polynomials of degree below $v$ vanish. Elimination by degree
proves surjectivity and that the full kernel is $\mathcal P_{v-1}$.
The quotient is represented isometrically by
$\phi_{v,s,c_0},\ldots,\phi_{D+v,s,c_0}$.
Set
\[
 w(t)=-i\arctan t=\tfrac12\log\frac{1-it}{1+it},\qquad
 g(t)=t^{-v}e^{-4w(t)}E_A(w(t)),\qquad
 g(0)=(-i)^v\mu_v/v! .                                  \tag{LT5}
\]
The zero-order definition of $v$ removes the singularity at zero.
Apply $\mathcal T_A$ to LT3 at $y=(S-c_0)/i$.
Since $c_0-c_1=4$, its image is $t^vg(t)$ times the same generating
function at $y=(S'-c_1)/i$. Comparing coefficients, including
$n!/\sqrt{h_{n,s}}=\rho_n/\sqrt{M_s}$, gives the exact matrix
\[
 A_{D,s}=\operatorname{diag}(\rho_i)_{i=0}^D^{-1}
             T_D(g)\operatorname{diag}(\rho_{i+v})_{i=0}^D,
 \quad T_D(g)_{ij}=\begin{cases}g_{j-i},&i\le j,\\0,&i>j.\end{cases}
                                                               \tag{LT6}
\]
Both copies of $M_s$ remain in the original norms; they cancel only
in this matrix connecting two frames with the same original order.
No scalar $v!/\mu_v$ is inserted.

\subsection{The finite analytic constants in LRC6}

Define, exactly as in the incoming calculation,
\[
\begin{gathered}
 g_\flat(t)=(1+t^2)^4g(t),\qquad
 F_0=4+4\sqrt{\delta^2+\gamma^2},\\
 B_0^{\flat}=2^{v+8}A_{\rm abs}e^{2\pi\gamma},\qquad
 B_2^{\flat}=B_0^{\flat}
 \{F_0(F_0+1)+4vF_0+4v(v+1)\}.
\end{gathered}                                                   \tag{LT7}
\]
The superscript on $B_0^{\flat},B_2^{\flat}$ only prevents collision
with WCF14's different zero-free-disk constant $B_0$; their numerical
values are the incoming LRC6 constants, without alteration.
For $\lambda_{uw}=(\beta_{8,u,w}-4)/2$, put
$\alpha_{uw}=4+\lambda_{uw}$ and $\beta^{\flat}_{uw}=4-\lambda_{uw}$.
Then
\[
 t^vg_\flat(t)=H_\flat(t):=
 \sum_{u,w}a_{uw}(1-it)^{\alpha_{uw}}(1+it)^{\beta^{\flat}_{uw}},
 \quad |\Re\lambda_{uw}|\le4\delta<2,
 \quad |\Im\lambda_{uw}|\le4\gamma,
 \quad |\alpha_{uw}|,|\beta^{\flat}_{uw}|\le F_0.            \tag{LT8}
\]
For $|t|<1$, $z=(1-it)/(1+it)$ satisfies
$\Re z=(1-|t|^2)/|1+it|^2>0$. Each $1\pm it$ has positive
real part, and the branches fixed at zero obey
$\log z=\log(1-it)-\log(1+it)$ throughout the disk.
For nonnegative integers $j,k$ with $j+k\le2$, the real exponents
in $(1-it)^{\alpha_{uw}-j}(1+it)^{\beta^{\flat}_{uw}-k}$ are
positive: the smallest possible one is $2-4\delta>0$.
Their sum is $8-j-k$, while the imaginary contribution to its
modulus is $\exp[-\Im\lambda_{uw}\arg z]\le e^{2\pi\gamma}$.
Consequently
\[
 |(1-it)^{\alpha_{uw}-j}(1+it)^{\beta^{\flat}_{uw}-k}|
 \le 2^{8-j-k}e^{2\pi\gamma}.                              \tag{LT9}
\]
For $f=(1-it)^\alpha(1+it)^\beta$ its first derivative has
coefficients $-i\alpha,i\beta$, and its second derivative has
coefficients $-\alpha(\alpha-1),2\alpha\beta,-\beta(\beta-1)$.
LT9 therefore proves
\[
\begin{split}
 |H_\flat|&\le2^8A_{\rm abs}e^{2\pi\gamma},\\
 |H_\flat'|&\le2^7(2F_0)A_{\rm abs}e^{2\pi\gamma},\\
 |H_\flat''|&\le2^6\{2F_0(F_0+1)+2F_0^2\}
                      A_{\rm abs}e^{2\pi\gamma}
 \le2^8F_0(F_0+1)A_{\rm abs}e^{2\pi\gamma}.
\end{split}                                                     \tag{LT10}
\]
On $1/2\le|t|<1$, use $g_\flat=t^{-v}H_\flat$ and
\[
 g_\flat''=t^{-v}H_\flat''-2vt^{-v-1}H_\flat'
                         +v(v+1)t^{-v-2}H_\flat .
\]
It follows that $|g_\flat|\le B_0^{\flat}$ and
$|g_\flat''|\le B_2^{\flat}$ on this annulus. Both functions
are holomorphic on the disk, since $H_\flat$ has a zero of order
$v$ at zero. Applying the maximum-modulus principle on each circle
of radius $r\in[1/2,1)$ proves
\[
 \boxed{|g_\flat(t)|\le B_0^{\flat},\qquad
 |g_\flat''(t)|\le B_2^{\flat}\quad(|t|<1).}              \tag{LRC6}
\]
If $g_\flat(t)=\sum_{n\ge0}b_nt^n$, Cauchy's integral formula
on $|t|=r<1$, followed by $r\uparrow1$, gives
\[
 |b_0|,|b_1|\le B_0^{\flat},\qquad
 |b_n|\le B_2^{\flat}/(n(n-1))\quad(n\ge2).               \tag{LT11}
\]
This argument has retained every interior zero of $g$.

\subsection{Exact matrix factorization and every exterior rank}

Let $S_D$ be the upper shift, with entries $1$ in positions $(i,i+1)$,
and set
\[
 R=\operatorname{diag}(\rho_i)_{i=0}^D,\qquad
 W=R^{-1}S_DR,\qquad V_v=\operatorname{diag}(\rho_{i+v}/\rho_i)_{i=0}^D,
 \qquad U_\flat=R^{-1}T_D(g_\flat)R .                    \tag{LT12}
\]
For any two series analytic at zero, $T_D(fh)=T_D(f)T_D(h)$:
the $(i,j)$ entry is exactly the finite convolution for coefficient
$j-i$, and all entries below the diagonal vanish. Since $S_D^{D+1}=0$,
all power-series evaluations at $S_D$ and at $W$ are finite.
The identity $g=g_\flat(1+t^2)^{-4}$ and LT6 thus give
\[
 \boxed{A_{D,s}=U_\flat(I+W^2)^{-4}V_v.}                 \tag{LRC7}
\]
This fixes the position of $V_v$ on the right. It is never commuted
through the other two matrices.

The nonzero images of different coordinate vectors under $W^n$ are
orthogonal, so its norm is the maximum of $\rho_{i+n}/\rho_i$.
At $s=a$, each factor $(i+j+1)/(i+j+s/2)$ is at most one, hence
$\|W^n\|\le1$ and $\|V_v\|\le1$. At $s=1$,
\[
 \frac{\rho_{i+n}^2}{\rho_i^2}
 =\prod_{j=0}^{n-1}\frac{i+j+1}{i+j+1/2}
 \le\frac{n!}{(1/2)_n}\le2n+1.                         \tag{LT13}
\]
The first inequality follows because every factor decreases with $i$.
For the second, induction starts at $n=0$; multiplying $2n+1$ by
$(n+1)/(n+1/2)$ gives $2n+2\le2n+3$.
In particular $\|V_v\|\le\sqrt{2v+1}$.
Now $U_\flat=\sum_{n=0}^D b_nW^n$. For $n\ge2$,
\[
 \frac{\sqrt{2n+1}}{n(n-1)}\le\sqrt{10}\,n^{-3/2},\qquad
 \sum_{n=2}^\infty n^{-3/2}
 \le\int_1^\infty x^{-3/2}\,dx=2.
\]
LT11--13 prove the original finite constant
\[
 \|U_\flat\|\le H:=B_0^{\flat}(1+\sqrt3)+2\sqrt{10}B_2^{\flat}.
                                                               \tag{LT14}
\]
Put
\[
 H_s=\begin{cases}\max\{1,H\sqrt{2v+1}\},&s=1,\\
                  \max\{1,H\},&s=a,\end{cases}
 \qquad
 \kappa_s=\begin{cases}1+\sqrt{8/3},&s=1,\\2,&s=a.\end{cases}
                                                               \tag{LT15}
\]
Writing $Q_D=I+W^2$, triangularity gives $\det Q_D=1$.
For $s=1$, the decreasing-factor argument gives the sharper
$\|W^2\|\le\rho_2/\rho_0=\sqrt{8/3}$, including $D<2$ when
$W^2=0$. For $s=a$ it gives $\|W^2\|\le1$.
Thus $\|Q_D\|\le\kappa_s$ in both cases.

For any invertible complex $n\times n$ matrix $T$ with singular
values $\sigma_1\ge\cdots\ge\sigma_n>0$, diagonalizing its
positive squared norm, or taking its singular-value decomposition,
gives
\[
 \|\textstyle\bigwedge^rT\|=\prod_{j=1}^r\sigma_j,
 \qquad
 \|\textstyle\bigwedge^rT^{-1}\|
 =\frac{\|\bigwedge^{n-r}T\|}{|\det T|}
 \quad(0\le r\le n).                                  \tag{LT16}
\]
The second equality multiplies the largest inverse singular values;
its denominator is the full determinant, not a truncated product.
Apply it to $Q_D^{\oplus c}$, where $c\in\{1,4\}$, whose dimension
is $cd$ and determinant one. Exterior functoriality and
submultiplicativity give
\[
 \|\bigwedge^r\bigl((Q_D^{-4})^{\oplus c}\bigr)\|
 \le\|\bigwedge^r\bigl((Q_D^{-1})^{\oplus c}\bigr)\|^4
 =\|\bigwedge^{cd-r}\bigl(Q_D^{\oplus c}\bigr)\|^4
 \le\kappa_s^{4(cd-r)}.                                 \tag{LT17}
\]
This is valid at $r=cd$, where the right side equals one.
For a complex matrix $A=U\Sigma V^*$,
$A^{\mathsf T}=\overline V\Sigma U^{\mathsf T}$ is another
singular-value decomposition; ordinary transposition therefore
preserves every exterior norm. Apply LRC7 to $c$ orthogonal copies,
use LT14 for the first factor and LT13 for the last, and then LT17.
For $1\le r\le cd$ this proves exactly
\[
 \boxed{\log\|\bigwedge^r\bigl((A_{D,s}^{\mathsf T})^{\oplus c}\bigr)\|
 \le F_{D,s,r,c}:=r\log H_s+4(cd-r)\log\kappa_s.}         \tag{LRC8}
\]
Set $F_{D,s,0,c}=0$, the exact logarithm of the empty exterior norm.

LT6 also proves the full complex determinant identity
\[
 \det A_{D,s}=\bigl((-i)^v\mu_v/v!\bigr)^d
                         \prod_{i=0}^{D}\frac{\rho_{i+v}}{\rho_i}.
                                                               \tag{LT18}
\]
The phase is retained in this equality. At $s=1$ the product of
positive ratios is at least one. At $s=a$,
\[
 \log\prod_{i=0}^{D}\frac{\rho_i}{\rho_{i+v}}
 =\tfrac12\sum_{i=0}^{D}\sum_{j=0}^{v-1}
       \log\left(1+\frac{s/2-1}{i+j+1}\right)
 \le\tfrac12(s/2-1)v\sum_{j=1}^d\frac1j .                \tag{LT19}
\]
Here the elementary inequality $\log(1+x)\le x$ applies to each
nonnegative $x$; the sums are empty if $v=0$.
Define $\log^+x=\max\{0,\log x\}$, $H_d^{\rm harm}=\sum_{j=1}^d1/j$,
and retain WCF's exact finite allowance
\[
 J_{D,s}=d\log^+(1/|g(0)|)
       +\tfrac12(s/2-1)_+vH_d^{\rm harm}.                \tag{LT20}
\]
LT18--19 give $|\det A_{D,s}|^{-1}\le e^{J_{D,s}}$.
Apply LT16 to the whole matrix
$B=(A_{D,s}^{\mathsf T})^{\oplus c}$ and LRC8 to its complementary
rank. For $1\le r\le cd$,
\[
 \boxed{\log\|\bigwedge^r\bigl((A_{D,s}^{-\mathsf T})^{\oplus c}\bigr)\|
 \le I_{D,s,r,c}:=cJ_{D,s}+F_{D,s,cd-r,c}.}              \tag{LRC9}
\]
Set $I_{D,s,0,c}=0$. This does not assert that the displayed
positive-rank formula extends to rank zero. At full rank it gives
$I_{D,s,cd,c}=cJ_{D,s}$; at $D=0$ the original matrix is the scalar
$g(0)\rho_v$. Thus the smallest cutoff, $v=0$, empty exterior powers
and full exterior powers are all included. No assumption on absence
or number of interior zeros of $g$ has entered these estimates.

There is an explicit constant for the claimed uniform order. Suppose
$D\le C_{\rm cut}a^2$ with fixed $C_{\rm cut}\ge0$, and put
\[
 \begin{gathered}
 K=C_{\rm cut}+1,\qquad G=\log^+(1/|g(0)|),\qquad
 H_* =\max\{1,H\sqrt{2v+1}\},\qquad
 \kappa_*=1+\sqrt{8/3},\\
 L_* =\max\{\log H_*,4\log\kappa_*\},\qquad
 J_*=KG+\tfrac v4(3+\log K).
 \end{gathered}                                                   \tag{LT21}
\]
Then $d\le Ka^2\le Kq$ and
$H_d^{\rm harm}\le1+\log d\le1+\log K+2\log a$.
Since $\log a\le a$, $a\le q$ and $a^2\le q$, LT20 implies
$J_{D,s}\le J_*q$. Moreover the expression in LRC8 is bounded by
$cdL_*$. Therefore for every retained rank, including zero,
\[
 0\le F_{D,s,r,c}\le cKL_*q,\qquad
 0\le I_{D,s,r,c}\le c(J_*+KL_*)q .                    \tag{LT22}
\]
All constants depend only on the actual conductor and the displayed
cutoff constant. These inequalities prove uniform $O(q)$ bounds at
all ranks and both original source orders, with no rank restriction
beyond $0\le r\le c(D+1)$.

\subsection{The literal lower covariance and every attained quotient}

Now impose the original endpoint domain $N\ge q-1$, $D=N-v$.
Let $\mathsf E_-$ and $\mathsf E_+$ be the full original observation
matrices in the unchanged root order:
\[
 (\mathsf E_-)_{n,(u,w)}=\phi_{n,s,c_1}(\beta_{a-8,u,w}),\ 0\le n\le D,
 \quad
 (\mathsf E_+)_{n,(u,w)}=\phi_{n,s,c_0}(\beta_{a,u,w}),\ 0\le n\le N.
                                                               \tag{LT23}
\]
The roots are distinct because their real and imaginary parts recover
the two indices. Since $q-Q=16a-48\ge96>v$, we have $D\ge Q-1$.
The leading coefficients in LT2 and the Vandermonde determinant prove
that $\mathsf E_-$ has rank $Q$ and $\mathsf E_+$ has rank $q$.
Define the original conductor matrix $C_a$ by collecting every term
in the identity
\[
 \sum_{u,w}(C_a^{\mathsf T}\eta)_{uw}e^{\beta_{a,u,w}z}
 =E_A(z)\sum_{u,w=0}^{a-8}\eta_{uw}e^{\beta_{a-8,u,w}z}.   \tag{LT24}
\]
This uses exactly
$\beta_{a-8,u,w}+\beta_{8,r,t}=\beta_{a,u+r,w+t}$.
All contributions producing the same root are retained.
If the left side is zero, the lower exponential sum vanishes on
an open set where $E_A\ne0$, and hence everywhere. Its derivative
Vandermonde forces $\eta=0$; thus $C_a^{\mathsf T}$ is injective.

For a lower functional $\Lambda$, complex linearity gives
$\Lambda(\mathcal T_A\phi_{v+j,s,c_0})
=\sum_i(A_{D,s})_{ij}\Lambda(\phi_{i,s,c_1})$.
This proves the ordinary-transpose square, including the vanishing
first $v$ rows,
\[
 \mathsf E_+C_a^{\mathsf T}
 =\begin{bmatrix}0_{v\times Q}\\A_{D,s}^{\mathsf T}\mathsf E_-\end{bmatrix},
 \quad K_{-,D,s}=\mathsf E_-^*\mathsf E_->0,\quad
 K_{C,N,s}=\mathsf E_-^*\overline{A_{D,s}}A_{D,s}^{\mathsf T}\mathsf E_->0.
                                                               \tag{LT25}
\]
This is the original WCF22 map with the unchanged orders $s=1,a$,
the lower centre $a/2-4$, and the shifted cutoff $N-v$.

For independent columns $e_1,\ldots,e_r$, the determinant of their
Gram matrix is $\|e_1\wedge\cdots\wedge e_r\|^2$. This follows
by expanding both sides as the determinant of their pairwise inner
products; it is the classical exterior Gram identity, for example
Lyons, \S4, (4.1) \cite{human:lyons2003}.
Apply $\bigwedge^r B$ to this wedge and $\bigwedge^r B^{-1}$
to the image. LRC8--9 give the two directed determinant bounds.
Taking $E=\mathsf E_-$, $B=A_{D,s}^{\mathsf T}$ and $r=Q$ proves
\[
 \boxed{-2I_{D,s,Q,1}\le e^C_{N,s}
 :=\log\det K_{C,N,s}-\log\det K_{-,N-v,s}
 \le2F_{D,s,Q,1}.}                                      \tag{LRC10}
\]

For completeness, the same argument controls every compatible attained
quotient, rather than only a restriction. Let $K_0\subset Y\subset H$
be subspaces of a finite-dimensional Hilbert space and $B:H\to H'$
an invertible linear map. Put
$Y_0=Y\cap K_0^\perp$ and $Y_1=BY\cap(BK_0)^\perp$.
Every affine fibre modulo $K_0$ has a unique minimum-norm element in
$Y_0$, by its orthogonal decomposition. The induced map and its inverse
in these minimum representatives are exactly
\[
 \overline B=P_{Y_1}B|_{Y_0},\qquad
 \overline B^{-1}=P_{Y_0}B^{-1}|_{Y_1}.                  \tag{LT26}
\]
Indeed subtraction in the first formula lies in $BK_0$ and subtraction
in the second lies in $K_0$, proving that the compositions are identities
on quotient classes and hence on their unique minimum representatives.
Projections and isometric inclusions contract every exterior norm,
as is immediate from their zero and unit singular values.
It follows that each exterior norm of $\overline B$ and its inverse
is bounded by the corresponding exterior norm of $B$ and $B^{-1}$.
For any common quotient coordinate frame of rank $r$, the quotient
Gram determinant ratio therefore lies in
\[
 -2I_{D,s,r,c}\ \le\ \log\frac{\det G_1}{\det G_0}
                    \ \le\ 2F_{D,s,r,c}                 \tag{LT27}
\]
when $B=(A_{D,s}^{\mathsf T})^{\oplus c}$.
A succession of restrictions and quotients reduces to the full
preimage pair $K_0\subset Y$ in the original ambient space: the
union of the affine fibres minimized at successive steps is exactly
the final affine fibre. Its minimum is attained by the same orthogonal
projection. Thus the complete denominators remain, and no extra factor
is paid for the dimension of a denominator or the number of steps.

Retain $N_i=(q-1,q,2q-1,2q)_i$, $D_i=N_i-v$ and
$\mathcal R f=-f(N_0)-f(N_1)+f(N_2)+f(N_3)$.
With $f_i^C=F_{D_i,s,Q,1}$ and $i_i^C=I_{D_i,s,Q,1}$, LRC10 gives
\[
 -2(f_0^C+f_1^C+i_2^C+i_3^C)
 \le\mathcal Re^C\le
 2(i_0^C+i_1^C+f_2^C+f_3^C).                            \tag{LT28}
\]
The actual cutoffs have $D_i\le2(a+1)^2\le3a^2$ for $a\ge9$;
LT22 proves that both finite allowances are $O_{\rm actual}(q)$.
The precise native $\Xi$ application of LT27, including its full
annihilators, coordinate section and original lower Gram denominators,
is given next. All earlier independently proved upper and lower
allowances may be retained by taking the minimum separately on each
directed side. The symbol, its zeros and their pole directions have
never been removed from the original map.
\endgroup


% Exact Xi receiver, reconstructed from NCT1--14,20 and TAC1b--i,5,7,9--11.
% The containing transport proof supplies the proved exterior bounds F and I.
\begingroup
\subsection{The exact analytic comparison and its four signed endpoints}
\label{sec:xi-exact-map}

Retain the actual simple quartet, original period and logarithm branch,
all coefficients of the original nonzero conductor, and the original root
order. Write
\[
\begin{gathered}
 a=4l+1\ge9,\quad q=(a+1)^2,\quad Q=(a-7)^2=q',\quad
 m=q-Q-v,\quad N\ge q-1,\quad D=N-v,\quad
 d=D+1,\quad t=N-q+1,\\
 s\in\{1,a\},\quad c_0=a/2,\quad c_1=a/2-4,\quad
 \beta_{n,u,w}=n/2+(2u-n)\delta+i(2w-n)\gamma,\\
 E_A(z)=\sum_{r,h=0}^{8}a_{rh}e^{\beta_{8,r,h}z},\quad
 \mu_j=E_A^{(j)}(0),\quad v=\operatorname{ord}_0E_A\le80,\quad
 \mu_v\ne0,\quad 0<\delta<1/2,\quad\gamma>2.
\end{gathered}\tag{XIM1}
\]
Thus $t\ge0$, $q-Q=16a-48\ge96$, and $m>0$. All polynomial
spaces have complex coefficients. The original Hermitian norm and frame are
\[
\begin{gathered}
 dm_s(y)=\frac{(2\pi)^{s/2}2^{s/2}}{4\pi\Gamma(s/2)}
           |\Gamma(s/4+iy/2)|^2\,dy,\qquad M_s=(2\pi)^{s/2},\\
 p_{0,s}=1,\quad p_{1,s}=y,\quad
 p_{n+1,s}=yp_{n,s}-n(n-1+s/2)p_{n-1,s},\\
 h_{n,s}=M_sn!(s/2)_n,\qquad
 \phi_{n,s,c}(S)=p_{n,s}((S-c)/i)/\sqrt{h_{n,s}},\qquad
 \|P\|_{s,c}^2=\int_{\mathbb R}|P(c+iy)|^2dm_s(y).
\end{gathered}\tag{XIM2}
\]
The frame is the original orthonormal frame, with the unchanged
Meixner--Pollaczek dictionary and norm
\cite[eqs.~(18.19.8)--(18.19.9), (18.22.8)]{human:dlmf18}; its leading coefficient is
$i^{-n}/\sqrt{h_{n,s}}$. Neither the mass nor this phase is changed.
The duals below are complex-linear duals, with the dual Hilbert norms.

Put $\mathcal P_j=\mathbb C[S]_{\le j}$, with $\mathcal P_{-1}=0$,
and use $S'$ for the target variable. The unscaled conductor map and its
matrix in the displayed orthonormal frames are
\[
\begin{gathered}
 \mathcal T_AP(S')=\sum_{r,h=0}^{8}a_{rh}P(S'+\beta_{8,r,h}),\\
 A=A_{D,s}:\mathcal P_N/\mathcal P_{v-1}\longrightarrow
             \mathbb C[S']_{\le D},\\
 \mathcal T_A\phi_{v+j,s,c_0}
     =\sum_{i=0}^D A_{ij}\phi_{i,s,c_1},\qquad 0\le j\le D,\\
 \rho_n=\sqrt{n!/(s/2)_n},\quad
 w(z)=-i\arctan z,\quad
 g(z)=z^{-v}e^{-4w(z)}E_A(w(z))=\sum_{j\ge0}g_jz^j,\\
 g_0=(-i)^v\mu_v/v!,\qquad
 A_{ij}=\begin{cases}g_{j-i}\rho_{j+v}/\rho_i,&i\le j,\\
                         0,&i>j.
          \end{cases}
\end{gathered}\tag{XIM3}
\]
These are the literal NCT3--5 and WCF4--6 matrices. Indeed
$\mathcal T_AS^n=\sum_{j=v}^n\binom nj\mu_jS'^{n-j}$, with
nonzero leading coefficient $\binom nv\mu_v$ when $n\ge v$.
It follows by elimination in degree that the kernel is exactly
$\mathcal P_{v-1}$ and that $A$ is an isomorphism. The quotient
norm uses its original orthogonal representatives
$\phi_{v,s,c_0},\ldots,\phi_{N,s,c_0}$.
For the matrix entries, multiplying the recurrence in XIM2 by
$z^n/n!$ gives the differential equation
$(1+z^2)\partial_zG=(y-(s/2)z)G$, $G(y,0)=1$, whose solution is
$G(y,z)=(1+z^2)^{-s/4}\exp(y\arctan z)$.
Applying $\mathcal T_A$ to this generating identity centred at $c_0$
multiplies the same identity centred at $c_1$ by $z^vg(z)$, since
$c_0-c_1=4$. Coefficient comparison and
$n!/\sqrt{h_{n,s}}=\rho_n/\sqrt{M_s}$ give exactly XIM3.

Here are the full relation and annihilator spaces. Define
\[
 \chi_a(S)=\prod_{u,w=0}^{a}(S-\beta_{a,u,w}),\qquad
 \chi_-(S')=\prod_{u,w=0}^{a-8}(S'-\beta_{a-8,u,w}).
 \tag{XIM4}
\]
Use brackets $[P]_-$ for the coefficient column of $P$ in
$\phi_{0,s,c_1},\ldots,\phi_{D,s,c_1}$, and $[P]_+$ for the
coefficient column of its class modulo $\mathcal P_{v-1}$ in
$\phi_{v,s,c_0},\ldots,\phi_{N,s,c_0}$. Form the matrices
\[
\begin{gathered}
 \mathsf X_N=\bigl[[\chi_aS^j]_+\bigr]_{0\le j<t}
       \in\mathbb C^{d\times t},\qquad
 \mathsf R_N=\bigl[[\mathcal T_A(\chi_aS^j)]_-\bigr]_{0\le j<t}
       =A\mathsf X_N\in\mathbb C^{d\times t},\\
 \mathsf L_N=\bigl[[\chi_-S'^j]_-\bigr]_{0\le j<d-Q}
       \in\mathbb C^{d\times(d-Q)}.
\end{gathered}\tag{XIM5}
\]
At $t=0$, the first two matrices have no columns. The upper relation
columns are independent modulo $\mathcal P_{v-1}$ because their
nonzero linear combinations have degree at least $q>v-1$. The
native relation columns have distinct degrees $q+j-v$ and nonzero
leading coefficients $\binom{q+j}{v}\mu_v$. The ideal columns are
monic of distinct degrees $Q+j$. Thus their respective ranks are
$t,t,d-Q$.

For every lower root, the identity
$\beta_{a-8,u,w}+\beta_{8,r,h}=\beta_{a,u+r,w+h}$ shows that
$\mathcal T_A(\chi_aS^j)$ vanishes at that root. Distinctness of
the roots follows by comparing real and imaginary parts. Division
by the monic $\chi_-$ therefore has zero remainder. Consequently
$\operatorname{im}\mathsf R_N\subset\operatorname{im}\mathsf L_N$.
For a functional with coordinate column
$z=(\Lambda(\phi_{i,s,c_1}))_{i=0}^D$, evaluation on a polynomial
with coefficient column $p$ is $z^{\mathsf T}p$. Its norm is
$\|z\|_2$. The two exact annihilator spaces are therefore
\[
\begin{gathered}
 \mathscr Y_D=\ker\mathsf R_N^{\mathsf T}
    =(\operatorname{im}\overline{\mathsf R_N})^\perp,
       \qquad \dim\mathscr Y_D=d-t=q-v,\\
 \mathscr K_D=\ker\mathsf L_N^{\mathsf T}
    =(\operatorname{im}\overline{\mathsf L_N})^\perp,
       \qquad \dim\mathscr K_D=Q,\\
 \mathscr K_D\subset\mathscr Y_D,\qquad
 \dim(\mathscr Y_D/\mathscr K_D)=m.
\end{gathered}\tag{XIM6}
\]
Complex conjugation in the orthogonal-complement expressions is
necessary because the algebraic annihilator uses ordinary transpose.

We identify these spaces with the original analytic functionals and
lower evaluations, including their frames. Let
$\mathbb V[n,\beta]=\beta^n$, $0\le n<q$, be the upper
Vandermonde in the original root order, and set
\[
\begin{gathered}
 \mathscr N_v=\{\nu\in\mathbb C^q:(\mathbb V\nu)_n=0
                                     \ (0\le n<v)\},\\
 C_{(u,w),(i,j)}=
 \begin{cases}a_{i-u,j-w},&0\le i-u,j-w\le8,\\0,&\text{otherwise},\end{cases}
       \qquad C\in\mathbb C^{Q\times q},\\
 \mathcal N_\nu(z)=\sum_\beta\nu_\beta e^{\beta z},\qquad
 \mathcal N_{C^{\mathsf T}\eta}(z)
   =E_A(z)\sum_{u,w=0}^{a-8}\eta_{uw}e^{\beta_{a-8,u,w}z}.
\end{gathered}\tag{XIM7}
\]
The last identity collects all original shift coefficients. Its zero
left side forces the lower exponential sum to vanish wherever
$E_A\ne0$, hence everywhere by analyticity; the lower Vandermonde
then forces $\eta=0$. Thus $C^{\mathsf T}$ is injective and its
image lies in $\mathscr N_v$.

Let $P_v$ insert the last $q-v$ coordinates into $\mathbb C^q$ and
put $J_v=P_v^{\mathsf T}\mathbb VC^{\mathsf T}$. Retain the original
first nonzero $Q$-row minor $\mathcal I$ of $J_v$, its complementary
row set $\mathcal J$, and the coordinate insertion $E_{\mathcal J}$.
Then
\[
 Z=\mathbb V^{-1}P_vE_{\mathcal J}\in\mathbb C^{q\times m},
 \qquad [C^{\mathsf T},Z]\text{ is a frame of }\mathscr N_v,
 \qquad F_b=\mathcal N_{Ze_b}/E_A.
 \tag{XIM8}
\]
The frame assertion follows by applying $\mathbb V$, taking the
$\mathcal I$ rows first, and then taking the complementary rows.
Each $F_b$ is analytic at zero because the first $v$ numerator
derivatives vanish.

These are also exactly the two packet annihilators in TAC1g--h.
Put $E=\mathbb C[S]/(\chi_a)$, $E_-=\mathbb C[S']/(\chi_-)$,
$L_v=\mathbb C[S]_{<v}\subset E$, and
$V_A=\ker(\overline{\mathcal T}_A:E\longrightarrow E_-)$.
Root addition proves that this packet map is well defined, and in
root-value coordinates it is $C$. Since $C^{\mathsf T}$ is injective,
the map is onto and $\dim V_A=q-Q$. The degree calculation in XIM3
gives $L_v\subset V_A$. The exact complex-linear pairing is
\[
\begin{gathered}
 \Phi:\mathbb C^q\longrightarrow E^\vee,\qquad
 \Phi(\nu)([P])=\sum_\beta\nu_\beta P(\beta),\\
 \Phi(\mathscr N_v)=\operatorname{ann}(L_v),\qquad
 \Phi(\operatorname{im}C^{\mathsf T})=\operatorname{ann}(V_A),\\
 \Psi:\mathscr N_v/\operatorname{im}C^{\mathsf T}
       \longrightarrow (V_A/L_v)^\vee,\quad
 \Psi([\nu])([P]+L_v)=\sum_\beta\nu_\beta P(\beta).
\end{gathered}\tag{XIM8a}
\]
The upper Vandermonde makes $\Phi$ invertible. Its first $v$ moment
equations are precisely vanishing on $L_v$. The conductor dual
functionals vanish on $V_A$, and both that image and
$\operatorname{ann}(V_A)$ have dimension $Q$, proving their equality.
Restriction therefore has exactly this kernel. Every functional on
$V_A/L_v$ extends linearly from $V_A$ to $E$ while vanishing on
$L_v$, proving that $\Psi$ is onto and hence an isomorphism.
At $v=0$, $L_v=0$ and the corresponding moment conditions are empty.

For $\nu\in\mathscr N_v$, let
$\Lambda_{\mathcal N_\nu/E_A}(S'^r)=(\mathcal N_\nu/E_A)^{(r)}(0)$.
Leibniz's finite identity proves
\[
 \Lambda_{\mathcal N_\nu/E_A}(\mathcal T_AP)
       =\Lambda_{\mathcal N_\nu}(P),\qquad \deg P\le N.
 \tag{XIM9}
\]
It vanishes for $P=\chi_aS^j$, so the native functional belongs
to $\mathscr Y_D$. If its first $D+1$ derivatives vanish, XIM9
makes all numerator derivatives through $N$ vanish. The first $q$
such equations have invertible Vandermonde, so $\nu=0$.
The dimension $q-v$ now proves that these are all of $\mathscr Y_D$.
Every lower evaluation annihilates the ideal generated by $\chi_-$.
Their first $Q$ evaluation rows form an invertible triangular change
of the lower Vandermonde, since $D\ge Q-1$. They are independent,
and dimension $Q$ proves that they are all of $\mathscr K_D$.

For completeness, the finite analytic observation matrix can be
constructed without an unevaluated division of germs. Put
$n_{jb}=\sum_\beta Z_{\beta b}\beta^j$ and $f_{rb}=F_b^{(r)}(0)$.
For $0\le r\le D$, successively set
\[
 f_{rb}=
 \frac{n_{v+r,b}-\displaystyle\sum_{j=v+1}^{v+r}
          \binom{v+r}{j}\mu_j f_{v+r-j,b}}
      {\binom{v+r}{v}\mu_v}.
 \tag{XIM10}
\]
At $r=0$ the sum is empty. All denominators are nonzero. This is
exactly the derivative product identity, hence retains every original
coefficient. Writing $\phi_{n,s,c_1}(S')=\sum_r\kappa_{nr}S'^r$,
define
\[
\begin{gathered}
 H=\mathsf E_-\in\mathbb C^{d\times Q},\qquad
 H_{n,\beta'}=\phi_{n,s,c_1}(\beta'),\\
 J\in\mathbb C^{d\times m},\qquad
 J_{nb}=\sum_{r=0}^n\kappa_{nr}f_{rb},\qquad
 \mathscr K_D=\operatorname{im}H,\quad
 \mathscr Y_D=\operatorname{im}[H,J].
\end{gathered}\tag{XIM11}
\]
Here $[H,J]$ has full column rank $Q+m=q-v$.

Let $\mathsf E_+[n,\beta]=\phi_{n,s,c_0}(\beta)$, $0\le n\le N$,
and let $\mathsf E_+^{\rm red}$ consist of rows $v,\ldots,N$.
The upper observation matrix has full column rank because its first
$q$ rows are an invertible triangular change of the upper Vandermonde.
The first $v$ rows vanish on $\mathscr N_v$. The ordinary dual
matrix $B=A^{\mathsf T}$ therefore gives the exact equations
\[
\begin{gathered}
 BH=\mathsf E_+^{\rm red}C^{\mathsf T},\qquad
 BJ=\mathsf E_+^{\rm red}Z,\qquad
 \mathsf E_+C^{\mathsf T}=\begin{bmatrix}0_{v\times Q}\\BH\end{bmatrix},
 \qquad \mathsf E_+Z=\begin{bmatrix}0_{v\times m}\\BJ\end{bmatrix},\\
 B\mathscr Y_D=\ker\mathsf X_N^{\mathsf T},\qquad
 B\mathscr K_D=\ker(A^{-1}\mathsf L_N)^{\mathsf T}
                =\operatorname{im}(\mathsf E_+^{\rm red}C^{\mathsf T}).
\end{gathered}\tag{XIM12}
\]
For the first line, apply each functional to the basis equation XIM3
and use XIM9. For the second line, use $\mathsf R_N=A\mathsf X_N$
and $B=A^{\mathsf T}$; invertibility proves both inclusions are
equalities. This proves the full commutative square on actual finite
jets. No adjoint and no scalar $\alpha_A=v!/\mu_v$ occurs.

The polynomial source of the analytic quotient is also exact.
Restriction of functionals induces isometric isomorphisms
\[
 \frac{\operatorname{ann}(\operatorname{im}\mathsf R_N)}
      {\operatorname{ann}(\operatorname{im}\mathsf L_N)}
   \simeq
 \left(\frac{\operatorname{im}\mathsf L_N}
               {\operatorname{im}\mathsf R_N}\right)^\vee,
 \quad
 \frac{\operatorname{ann}(\operatorname{im}\mathsf X_N)}
      {\operatorname{ann}(\operatorname{im}A^{-1}\mathsf L_N)}
   \simeq
 \left(\frac{\operatorname{im}A^{-1}\mathsf L_N}
               {\operatorname{im}\mathsf X_N}\right)^\vee.
 \tag{XIM13}
\]
The primal metrics first restrict the original polynomial metric and
then minimize over the full relation space. To prove the dual metric
claim, extend a functional from the larger displayed polynomial
subspace by composition with its orthogonal projection. This is the
minimum extension, since every other extension differs on the
orthogonal complement and the squared norms add. It still vanishes
on the relation subspace. All such extensions differ by precisely
the displayed denominator annihilator. This proves isometry as well
as the algebraic quotient assertion. The primal map is $A$ on these
two polynomial subquotients, and XIM12 is its contravariant ordinary
dual. Each primal subquotient has dimension $(d-Q)-t=m$.
The original packet identification is
\[
\begin{gathered}
 V_A/L_v\longrightarrow
 \operatorname{im}(A^{-1}\mathsf L_N)/\operatorname{im}\mathsf X_N,
 \quad [P]+L_v\longmapsto[[P]_+],\\
 Y_N=\mathbb C[S']_{\le D-Q}\big/
   \{\mathcal T_A(\chi_a h)/\chi_-:h\in\mathcal P_{t-1}\},\\
 \theta_N:V_A/L_v\longrightarrow Y_N,\qquad
 [P]+L_v\longmapsto[\mathcal T_AP/\chi_-].
\end{gathered}
 \tag{XIM13a}
\]
for the upper polynomial quotient and its divided native counterpart,
respectively. A representative of degree at most $N$ exists since
$N\ge q-1$. Changing it by $\chi_a h$ gives exactly an XIM5 relation,
and changing it by a polynomial of degree below $v$ gives zero.
Conversely, every polynomial in $\operatorname{im}(A^{-1}\mathsf L_N)$
has a degree-at-most-$N$ representative whose conductor vanishes on
the lower roots, hence represents $V_A$. A zero quotient class differs
from a multiple of $\chi_a$ by a polynomial of degree below $v$.
This proves both surjectivity and injectivity. Dividing the native
relation by $\chi_-$ is legitimate by XIM5; its relation map is the
original $h\mapsto\mathcal T_A(\chi_a h)/\chi_-$. Thus XIM13a is
exactly the TAC1d map, including its full relation space and scalar one.

We now give both attained coefficient minima. Put
\[
\begin{gathered}
 K_0=H^*H>0,\quad K_1=H^*B^*BH>0,\qquad
 P_0=HK_0^{-1}H^*,\quad P_1=BHK_1^{-1}H^*B^*,\\
 W_0=(I-P_0)J,\qquad W_1=(I-P_1)BJ,\\
 D_{a,D}^{(s)}=W_0^*W_0,\qquad
 \widehat D_{a,D}^{(s)}=W_1^*W_1.
\end{gathered}\tag{XIM14}
\]
The first metric is the original NCT10 native analytic metric. The
second is precisely NCT8/TAC5: setting $K_+=\mathsf E_+^*\mathsf E_+$,
XIM12 gives
\[
 \widehat D_{a,D}^{(s)}=
 Z^*\left[K_+-K_+C^{\mathsf T}
    (\overline C K_+ C^{\mathsf T})^{-1}\overline C K_+\right]Z.
 \tag{XIM15}
\]
For every $x\in\mathbb C^m$, completing the two positive squares
in the Schur-complement calculation
\cite[Appendix A.5.5, p.~650; Appendix C.4.1, p.~673]{human:boyd-vandenberghe}
proves
\[
\begin{aligned}
 x^*D_{a,D}^{(s)}x
   &=\min_{\eta\in\mathbb C^Q}\|Jx+H\eta\|_2^2,
 &\eta_0(x)&=-K_0^{-1}H^*Jx,\\
 x^*\widehat D_{a,D}^{(s)}x
   &=\min_{\eta\in\mathbb C^Q}\|B(Jx+H\eta)\|_2^2,
 &\eta_1(x)&=-K_1^{-1}H^*B^*BJx.
\end{aligned}\tag{XIM16}
\]
Both minima retain every lower evaluation and all their cross terms.
Each minimizer is unique because $K_0,K_1$ are positive definite.
More explicitly, the difference between the first squared norm and
its displayed minimum is
$(\eta-\eta_0(x))^*K_0(\eta-\eta_0(x))$, and the difference for
the second is $(\eta-\eta_1(x))^*K_1(\eta-\eta_1(x))$.
These identities follow on expanding, using the two displayed
formulas for $\eta_0,\eta_1$; each cross term cancels exactly.
If either minimum were zero, invertibility of $B$ and independence
of $[H,J]$ would imply $x=0$. Thus both analytic metrics are positive
definite. The two minimizing coefficient columns need not coincide.

The induced map on the quotient has the following exact domain,
codomain, and inverse:
\[
\begin{gathered}
 Q_0=\mathscr Y_D\cap\mathscr K_D^\perp=\operatorname{im}W_0,
 \quad Q_1=B\mathscr Y_D\cap(B\mathscr K_D)^\perp
                         =\operatorname{im}W_1,\\
 \overline B:Q_0\longrightarrow Q_1,
       \quad \overline B=P_{Q_1}B|_{Q_0},\qquad
 \overline B^{-1}=P_{Q_0}B^{-1}|_{Q_1},\qquad
 \overline B(W_0x)=W_1x.
\end{gathered}\tag{XIM17}
\]
The quotient identifications are isometric because each representative
has a unique orthogonal minimum lift. Projection after $B$ subtracts
an element of $B\mathscr K_D$; applying $B^{-1}$ turns that into an
element of $\mathscr K_D$, which the second projection removes.
This proves both inverse identities. It also proves the last equation
using XIM14. In the orthonormal minimum-lift frames
$W_0(D_{a,D}^{(s)})^{-1/2}$ and
$W_1(\widehat D_{a,D}^{(s)})^{-1/2}$, the matrix is
\[
 [\overline B]=
 (\widehat D_{a,D}^{(s)})^{1/2}(D_{a,D}^{(s)})^{-1/2},\qquad
 [\overline B]^*[\overline B]=
 (D_{a,D}^{(s)})^{-1/2}\widehat D_{a,D}^{(s)}(D_{a,D}^{(s)})^{-1/2}.
 \tag{XIM18}
\]
Indeed $W_1^*BW_0=W_1^*W_1=\widehat D_{a,D}^{(s)}$, because
the discarded part lies in $B\mathscr K_D$. Thus the original TAC9
correction is exactly
\[
 \epsilon_{a,N,s}
 =\log\det\widehat D_{a,D}^{(s)}-\log\det D_{a,D}^{(s)}
 =2\log|\det[\overline B]|.
 \tag{XIM19}
\]

Write the established all-rank logarithmic exterior estimates for the
actual ambient matrix, including its original weights, as
\[
 \log\left\|\bigwedge^r\bigl((A_{D,s}^{\mathsf T})^{\oplus c}\bigr)\right\|
       \le F_{D,s,r,c},\qquad
 \log\left\|\bigwedge^r\big((A_{D,s}^{\mathsf T})^{\oplus c}\big)^{-1}
       \right\|\le I_{D,s,r,c},\qquad 0\le r\le cd.
 \tag{XIM20}
\]
The $F$ and $I$ in this display are logarithms of exterior norms,
without the Gram factor two. Orthogonal projections contract every
exterior power: in orthonormal bases their singular values are zero
or one. Restrictions and inclusions of Hilbert subspaces also
contract exterior powers. Applying XIM17 in both directions and
then XIM20 at $c=1$, $r=m$, proves the finite bound
\[
 \boxed{-2I_{D,s,m,1}\le\epsilon_{a,N,s}\le2F_{D,s,m,1}.}
 \tag{XIM21}
\]
There is no added $Q$-rank allowance: the complete denominator has
already been retained in the exact quotient map and in both minima.
For $c$ copies, the denominator has dimension $cQ$ and the full
analytic quotient dimension $cm$. Any further restriction followed
by an attained quotient has its full preimages $K'\subset Y'$ in
the ambient $c$ copies; necessarily $K'$ contains the original
$cQ$-dimensional denominator. Nested minimization over these fibres
equals minimization over $K'$, because the union of the original
affine fibres is the final affine fibre. XIM17 applies to $K'\subset Y'$
and proves the same determinant interval $[-2I_{D,s,r,c},2F_{D,s,r,c}]$
with its actual final rank $r$, including $r=0$ with empty determinant
one. No intermediate dimension is charged a second time.

Finally retain the actual TAC endpoints and dual sign orientation:
\[
 N_i=(q-1,q,2q-1,2q)_i,\quad D_i=N_i-v,\quad
 \Xi_{a,s}=-\epsilon_{a,N_0,s}-\epsilon_{a,N_1,s}
             +\epsilon_{a,N_2,s}+\epsilon_{a,N_3,s}.
 \tag{XIM22}
\]
For these four endpoints,
$d_i=(q-v,q+1-v,2q-v,2q+1-v)_i$ and
$t_i=(0,1,q,q+1)_i$. The annihilator dimensions remain exactly
$q-v$ and $Q$, and the analytic quotient rank remains $m$ at every
endpoint. Set $F_i=F_{D_i,s,m,1}$ and $I_i=I_{D_i,s,m,1}$.
Negating the two low-endpoint inequalities in XIM21 and adding the
two high-endpoint inequalities gives
\[
 \boxed{
 -2(F_0+F_1+I_2+I_3)
       \le\Xi_{a,s}\le
  2(I_0+I_1+F_2+F_3).}
 \tag{XIM23}
\]
These are bounds on the same four correlated matrices appearing in
TAC9. A separate radius infimum, minimum of established allowances,
or improved ambient estimate can be inserted in each $F_i$ or $I_i$
provided it bounds XIM20; no attainment of an infimum is used.
The source order remains $1$ or $a$ at the lower grid, its centre
remains $a/2-4$, and its cutoff remains $N_i-v$. In particular no
replacement by the order $a-8$ or by an unshifted cutoff occurs.
The explicit uniform bounds LT21--22 of the preceding transport proof
apply because $D_i\le3a^2$ for $a\ge9$. Inserting them in XIM23
proves $\Xi_{a,s}=O_{\rm actual}(q)$ for both original source orders.
The identical endpoint bound XIM21 holds at every cutoff
$q-1\le N\le2q$, including each growing-prefix endpoint $N=q+R$
with $0\le R<q$, with the same fixed constant.
\endgroup
